\section{总结与延伸}

\subsection{整讲框架回顾}
这节课串起了三条线：评测线（什么指标值得优化）、系统线（harness 与工具如何设计）、安全线（为什么必须动态验证）。三条线最终汇聚到同一结论：\textbf{Agent 的竞争力来自系统闭环质量，而非单次文本能力}。

\begin{table}[H]
\centering
\begin{tabular}{p{3.1cm}p{4.4cm}p{4.2cm}p{2.6cm}}
\toprule
\textbf{主题} & \textbf{本讲结论} & \textbf{可落地做法} & \textbf{后续挑战} \\
\midrule
评测设计 & 评测驱动研发方向，且有寿命 & 持续更新任务分布，校准真实场景 & 抵抗泄露与过拟合 \\
Agent 架构 & ReAct + 工具反馈是基础单元 & 控制流分层设计，失败轨迹必审计 & 成本与稳定性平衡 \\
Harness 工程 & 工具接口需与模型能力共设计 & ACI 思维、紧凑动作、可压缩输出 & 通用性与专用性折中 \\
安全任务 & 必须从静态判断转向动态验证 & sanitizer/debugger 进入主循环 & 跨项目泛化与误报控制 \\
Big Sleep 实践 & 真实漏洞发现已可达成 & 以触发证据为目标函数做迭代 & 扩展到更复杂系统栈 \\
\bottomrule
\end{tabular}
\caption{Lecture 08 的核心结论总表}
\end{table}

\subsection{Further Reading}
\begin{itemize}
    \item Sutton 讲座中涉及的核心基线：MBPP、HumanEval、SWE-Bench、SWE-agent、AutoCodeRover。
    \item 安全评测与数据：公开漏洞数据库（CVE 生态）、CTF 基准、Meta 安全评测方向、DARPA 相关竞赛。
    \item 工程化方向：Agent Computer Interface、轨迹级失败分析、动态验证优先的 harness 设计。
    \item 安全研究实践：Big Sleep / Project Naptime 公开分享材料，以及与 Project Zero 协作的漏洞披露流程。
\end{itemize}

\subsection{本章小结}
Sutton 在收尾阶段的判断可以直接作为路线图：\textit{``This is a wide open area. We're just getting started.''} 对研究者和工程团队而言，下一步重点不是再做一轮表面 benchmark 刷分，而是把 Agent 放进更真实、更可审计、更可复现的安全工作流。

\end{document}
